Many cryptographic algorithms are considered unbreakable, regardless of the amount of computing power available.
These secure cryptosystems however do not run in a vacuum with spherical cows, and have to be implemented, either in
hardware or in software.\\
Regardless of the means of the implementation, the electronic devices that run these cryptosystems often
\textit{leak} information about the internal state of the device, potentially exposing secrets.
But how does a device leak information?

A \textit{side channel} is some observable aspect of a cryptosystem that reveals information held within that system~\cite{standaert2010introduction};
example of side channels can be the execution time of an algorithm, the power consumption of a device or the electromagnetic
radiation emitted while running some operations.\\
In this chapter we will describe the most common side channel attacks, how they can be used to retrieve secret information and
how to mitigate them.

\section{Timing attacks}
In a timing attack, the attacker tries to infer information about a secret by measuring the time it takes to execute some operation.

To better explain how timing attacks work, let us consider the following code snippet:
\begin{lstlisting}[language=C]
int check_pin(int *pin) {
    int correct_pin[4] = {3, 1, 4, 2};

    for (int i = 0; i < 4; i++) {
        if (pin[i] != correct_pin[i]) {
            return 0;
        }
    }
    return 1;
}
\end{lstlisting}
We can think this function as a part of a bigger system, where the user enters a PIN and the system checks if it is correct.
The function simply checks if the PIN entered by the user is equal to the correct PIN, and returns 1 if it is correct and 0 otherwise.
Being the PIN composed of decimal digits, the total amount of combinations is $10^4$, so on average we would need 5000 guesses to find
the correct PIN; however, we can reduce the number of guesses needed to find the correct PIN by using a timing attack.

We can observe that the function returns as soon as a digit of the PIN is incorrect: under the reasonable assumption that we can measure 
the execution time of the function, we can guess the correct PIN one digit at a time.
Starting from 0000, we cycle through all the possibilities for the first digit $(0000, 1000, \cdots, 9000)$ and measure the execution time:
out of the 10 measurements, the one with the highest execution time corresponds to the correct digit; we can then repeat the same process
for the remaining digits, so that after at most 40 guesses we find the correct pin.

Even though this attack strategy may seem trivial and easy to mitigate, it can still wreck havoc: real-life examples of timing attacks include, 
for example, Spectre~\cite{Kocher2018spectre} and Meltdown~\cite{Lipp2018meltdown}, which affected almost all modern processors and allowed
an attacker to read the memory of any process running on the same machine.

\section{Power attacks}
